%!TEX root = ../thesis.tex

%jama -- agent-based micro model chapter

\chapter{Agent-Based Micro Model}
\label{ch:abm}

\section{Overview and Design Rationale}
\label{sec:abm_overview}

%jama
Compartmental ordinary differential equation (ODE) models describe epidemic dynamics at the population level by averaging over all individuals.
This makes them analytically tractable and suitable for controller design, but it also limits their ability to represent spatial structure, heterogeneity in contact patterns, and the stochastic character of person-to-person transmission.

Agent-based models (ABMs) address these limitations by simulating each individual explicitly.
In the present work, each individual is represented as an agent with its own disease state, city assignment, and state-history counter.
Population-level epidemic behaviour then emerges from local interactions and stochastic transition rules rather than being imposed directly at the aggregate level \citep{zhang2025,polcz2025}.

Within the thesis, the micro model serves two purposes.
First, it acts as a simulation environment that converts prescribed transmission-rate trajectories into epidemic outcomes with agent-level variability.
Second, it provides aggregated compartment counts that can be compared with, or coupled back to, the macro-level epidemic model.

The implementation described here is organised around five MATLAB functions:
\texttt{init\_agents.m}, \texttt{update\_agents.m}, \texttt{aggregate\_stats.m}, \texttt{run\_micro\_sim.m}, and \texttt{run\_micro\_sim\_visual.m}.

\section{Agent States and Attributes}
\label{sec:agent_states}

\subsection{Compartment Structure}
\label{sec:compartments}

%jama
Each agent occupies one of seven mutually exclusive disease states at every time step.
The states mirror the compartment structure used by the macro-level model, which supports direct comparison between the two modelling scales.

\begin{table}[ht]
\centering
\caption{Agent state codes used in the micro model.}
\label{tab:states}
\begin{tabular}{@{}cll@{}}
\toprule
\textbf{Code} & \textbf{Symbol} & \textbf{Description} \\
\midrule
1 & $S$   & Susceptible \\
2 & $I$   & Infected (symptomatic, unvaccinated) \\
3 & $I_v$ & Infected (breakthrough, vaccinated) \\
4 & $R$   & Recovered \\
5 & $V$   & Vaccinated susceptible \\
6 & $H$   & Hospitalised \\
7 & $D$   & Dead \\
\bottomrule
\end{tabular}
\end{table}

\subsection{Agent Attributes}
\label{sec:attributes}

%jama
In addition to the disease state, each agent stores two further attributes:
\begin{itemize}
  \item \texttt{city}: an integer label (1 or 2) identifying the city to which the agent currently belongs.
  \item \texttt{days\_in\_state}: a non-negative integer counter that records how long the agent has remained in its current state.
    This variable supports future extensions with duration-dependent transition rules, even though the present implementation applies fixed daily probabilities.
\end{itemize}

\section{Population Initialisation}
\label{sec:init}

%jama
The function \texttt{init\_agents.m} creates an initial population of $N = N_1 + N_2$ agents, where $N_1$ and $N_2$ denote the populations of City~1 and City~2.
In the baseline configuration, $N_1 = N_2 = 500$, giving $N = 1\,000$.

All agents are first assigned to the susceptible state.
A randomly selected subset of $I_{0,1}$ agents in City~1 and $I_{0,2}$ agents in City~2 is then moved to the infected state to seed the outbreak.
Optional vaccinated cohorts $V_{0,1}$ and $V_{0,2}$ may also be introduced.
The default configuration is summarised in Table~\ref{tab:init_params}.

\begin{table}[ht]
\centering
\caption{Default initialisation parameters used in the micro model.}
\label{tab:init_params}
\begin{tabular}{@{}lll@{}}
\toprule
\textbf{Parameter} & \textbf{Description} & \textbf{Value} \\
\midrule
$N_1$ & Agents in City 1 & 500 \\
$N_2$ & Agents in City 2 & 500 \\
$I_{0,1}$ & Initial infected in City 1 & 10 \\
$I_{0,2}$ & Initial infected in City 2 & 8 \\
$V_{0,1}$ & Initial vaccinated in City 1 & 0 \\
$V_{0,2}$ & Initial vaccinated in City 2 & 0 \\
\bottomrule
\end{tabular}
\end{table}

A fixed random seed, \texttt{rng(1)}, is used in the simulation driver to improve reproducibility across runs.

\section{Disease Transmission Model}
\label{sec:transmission}

\subsection{Mean-Field Infection Probability}
\label{sec:mean_field}

%jama
The non-visual simulation uses a mean-field contact model that retains the same high-level infection logic as the macro formulation.
At time step $t$, the total infectious agents in city $c \in \{1, 2\}$ is
\begin{equation}
  \tilde{I}_c(t) = I_c(t) + I_{v,c}(t),
  \label{eq:total_infected}
\end{equation}
where $I_c(t)$ and $I_{v,c}(t)$ are the numbers of agents in states $I$ and $I_v$ in city $c$.

For a susceptible agent in city $c$, the daily infection probability is
\begin{equation}
  p_{\mathrm{inf}}^{(S)}(t) = \beta_c(t)\,\frac{\tilde{I}_c(t)}{N_c^{\mathrm{alive}}(t)},
  \label{eq:pinf_S}
\end{equation}
where $\beta_c(t)$ is the city-specific transmission rate and $N_c^{\mathrm{alive}}(t)$ denotes the number of non-dead agents in that city.
This corresponds to the frequency-dependent incidence structure used in compartmental epidemic models \citep{zhang2025}.

Vaccinated agents may also become infected, but with a reduced effective transmission rate:
\begin{equation}
  p_{\mathrm{inf}}^{(V)}(t) = \eta\,\beta_c(t)\,\frac{\tilde{I}_c(t)}{N_c^{\mathrm{alive}}(t)},
  \label{eq:pinf_V}
\end{equation}
where $\eta \in [0, 1]$ represents the vaccine-related reduction in infection risk.
In the present parameterisation, $\eta = 0.6$.

\subsection{Proximity-Based Infection in the Visual Simulation}
\label{sec:proximity}

%jama
The visual simulation replaces the mean-field assumption with a spatially explicit proximity rule.
Infectious agents can transmit disease to susceptible or vaccinated neighbours located within an interaction radius $r_{\mathrm{inf}} = 2.8$ spatial units and in the same city.

For an infectious agent $i$ and a candidate neighbour $j$ at distance $d_{ij}$, the transmission probability is modelled as
\begin{equation}
  p_{\mathrm{inf}}(i \to j,\, t) = \alpha\,\beta_c(t)\,\max\!\left(0,\, 1 - \frac{d_{ij}}{r_{\mathrm{inf}}}\right)\,\kappa_j,
  \label{eq:prox_inf}
\end{equation}
where $\alpha = 0.35$ is a contact-intensity scaling factor and
\begin{equation}
  \kappa_j = \begin{cases} 0.35, & \text{if agent } j \text{ is vaccinated}, \\ 1, & \text{otherwise.} \end{cases}
\end{equation}
The probability is clamped to $[0, 1]$ before the Bernoulli draw is evaluated.
This spatial variant reveals local clustering and mobility effects that are hidden in the aggregate mean-field formulation.

\section{Disease Progression Model}
\label{sec:progression}

%jama
State transitions other than infection are governed by independent daily Bernoulli probabilities.
Table~\ref{tab:transitions} summarises the transition rules used in the micro model.

\begin{table}[ht]
\centering
\caption{Daily transition probabilities in the micro model.}
\label{tab:transitions}
\begin{tabular}{@{}lllll@{}}
\toprule
\textbf{From} & \textbf{To} & \textbf{Parameter} & \textbf{Value} \\
\midrule
$S$ & $V$  & $\nu_c$ (vaccination rate in city $c$) & 0.00 (baseline) \\
$S$ & $I$  & $p_{\mathrm{inf}}^{(S)}$ from Eq.~\eqref{eq:pinf_S} & --- \\
\midrule
$V$ & $S$  & $\omega_V$ (waning immunity) & $1/140$ \\
$V$ & $I_v$& $p_{\mathrm{inf}}^{(V)}$ from Eq.~\eqref{eq:pinf_V} & --- \\
\midrule
$I$ & $H$  & $h$ (hospitalisation) & 0.05 \\
$I$ & $D$  & $\mu$ (fatality) & 0.001 \\
$I$ & $R$  & $\gamma$ (recovery) & $1/12$ \\
\midrule
$I_v$ & $H$ & $h_v$ (breakthrough hospitalisation) & 0.02 \\
$I_v$ & $D$ & $\mu_v$ (breakthrough fatality) & 0.0005 \\
$I_v$ & $R$ & $\gamma_v$ (breakthrough recovery) & $1/9$ \\
\midrule
$R$ & $S$  & $\omega_R$ (waning immunity after recovery) & $1/110$ \\
\midrule
$H$ & $D$  & $\mu_h$ (hospital fatality) & 0.02 \\
$H$ & $R$  & $\gamma_h$ (hospital discharge) & $1/25$ \\
\bottomrule
\end{tabular}
\end{table}

Transitions are evaluated sequentially within each day.
The first transition whose Bernoulli draw succeeds determines the next state for that agent, and subsequent transitions are skipped for that time step.
For example, an agent in state $I$ first faces hospitalisation, then death, and finally recovery if the previous events do not occur.

Both vaccinated and recovered agents may return to the susceptible state as immunity wanes.
Breakthrough infections are parameterised to have reduced severity compared with unvaccinated infections, reflected in lower hospitalisation and fatality probabilities and a shorter expected recovery time \citep{polcz2025}.

\section{Two-City Spatial Structure}
\label{sec:two_city}

%jama
The micro model represents a two-city metapopulation.
Each agent is associated with one of two cities, and the city label determines which transmission-rate input $\beta_c(t)$ is applied.
This structure makes it possible to analyse regional heterogeneity and city-specific intervention scenarios.

In the visual simulation, agents move within bounded spatial regions representing the two cities, separated by a corridor.
Inter-city commuting is modelled stochastically: each mobile agent may travel to the other city with probability $p_{\mathrm{commute}} = 0.015$ per day and remain there for a random stay of two to five days.
This extension captures how movement between regions can contribute to cross-city transmission \citep{zhang2025}.

In the non-visual simulation, agents remain assigned to a fixed city and mixing within each city is mean-field.
This version is computationally simpler and therefore more suitable for repeated simulation and controller coupling.

\section{Aggregation Interface to the Macro Model}
\label{sec:aggregation}

%jama
The function \texttt{aggregate\_stats.m} maps individual agent states to aggregate compartment counts.
At the end of each simulated day, the model computes:
\begin{align*}
  S(t) &= \#\{i : \mathrm{state}_i(t) = 1\}, &
  I(t) &= \#\{i : \mathrm{state}_i(t) = 2\}, \\
  I_v(t) &= \#\{i : \mathrm{state}_i(t) = 3\}, &
  R(t) &= \#\{i : \mathrm{state}_i(t) = 4\}, \\
  V(t) &= \#\{i : \mathrm{state}_i(t) = 5\}, &
  H(t) &= \#\{i : \mathrm{state}_i(t) = 6\}, \\
  D(t) &= \#\{i : \mathrm{state}_i(t) = 7\}. &
\end{align*}
These counts are returned globally and, where needed, separately for each city.
The resulting aggregated trajectories provide the interface between the individual-based simulation and any macro-level analysis or control layer.

Population is conserved throughout the simulation because each transition moves an agent from exactly one state to exactly one other state.
No births, immigration, or departures are included in the present model.
